\subsection{INVARIANT DIFFERENTIAL FORMS}

Let X be a locally finite dimensional real manifold of class \(C^{\infty}\) , and let \((g, x) \mapsto gx\) be a law of left operation of class \(C^{\infty}\) of a connected compact Lie group G on X. For \(g \in G\) , denote by \(\tau_{g}\) the automorphism \(x \mapsto gx\) of X. Denote by \(\Omega(X)\) the algebra of real differential forms of class \(C^{\infty}\) on X (Differentiable and Analytic Manifolds, Results, 8.3.1).

For any element \(\xi\) of \(\mathfrak{g}\) , denote by \(D_{\xi}\) the corresponding vector field on X (Chap. III, §3, no. 5) and by \(\theta(\xi), i(\xi)\) the corresponding operators on \(\Omega(X)\) , so that we have the formulas (Differentiable and Analytic Manifolds, Results, 8.4.5 and 8.4.7)

\[
\theta (\xi) \omega = d (i (\xi) \omega) + i (\xi) d \omega\tag{15}
\]

\[
\frac {d}{d t} (\tau_ {\mathrm{exp} t \xi} ^ {*} \omega) = \tau_ {\mathrm{exp} t \xi} ^ {*} (\theta (\xi) \omega).\tag{16}
\]

A differential form \(\omega \in \Omega(\mathrm{X})\) is invariant if \(\tau_g^*\omega = \omega\) for all \(g \in \mathrm{G}\) ; by formula (16), this is equivalent to \(\theta(\xi)\omega = 0\) for all \(\xi \in \mathfrak{g}\) . Denote by \(\Omega(\mathrm{X})^{\mathrm{G}}\) the space of invariant differential forms on X; if \(\omega \in \Omega(\mathrm{X})^{\mathrm{G}}\) , we have \(d\omega \in \Omega(\mathrm{X})^{\mathrm{G}}\) , so \(\Omega(\mathrm{X})^{\mathrm{G}}\) is a subcomplex of the complex \((\Omega(\mathrm{X}), d)\) .

THEOREM 2. The canonical injection \(\iota : \Omega(\mathrm{X})^{\mathrm{G}} \to \Omega(\mathrm{X})\) is a homotopism of complexes (Algèbre, Chap. X, p.~33, déf. 5); the map \(\omega \mapsto \omega^{\sharp} = \int_{\mathrm{G}} \tau_g^* \omega dg\) is a homotopism, inverse to it up to homotopy. In particular, the map \(\mathrm{H}(\iota) : \mathrm{H}(\Omega(\mathrm{X})^{\mathrm{G}}) \to \mathrm{H}(\Omega(\mathrm{X}))\) is bijective.

By Cor. 3 of no. 4, the map \(\omega \mapsto \omega^{\sharp}\) is a morphism of complexes from \(\Omega(\mathrm{X})\) to \(\Omega(\mathrm{X})^{\mathrm{G}}\) that induces the identity on the subcomplex \(\Omega(\mathrm{X})^{\mathrm{G}}\) ; thus, to prove the theorem it suffices to construct a homomorphism \(s: \Omega(\mathrm{X}) \to \Omega(\mathrm{X})\) , graded of degree \(-1\) , such that

\[
\omega^ {\sharp} - \omega = (d \circ s + s \circ d) (\omega) \quad \text { for   all } \omega \in \Omega (\mathrm{X}).\tag{17}
\]

By Lemma 1 of Integration, Chapter IX, §2, no. 4 and Remark 1 of §2, no. 2, there exists a positive measure \(d\xi\) on \(\mathfrak{g}\) of compact support whose image under the exponential map is equal to \(dg\) . For \(\omega \in \Omega(X)\) , put

\[
s (\omega) = \int_ {0} ^ {1} \left\{\int_ {\mathfrak {g}} \tau_ {\exp t \xi} ^ {*} (i (\xi) \omega). d \xi \right\} d t;
\]

we have to show that formula (17) is satisfied. As in the proof of Cor. 1 (no. 4), we verify the formula

\[
d s (\omega) = \int_ {0} ^ {1} \left\{\int_ {\mathfrak {g}} \tau_ {\exp t \xi} ^ {*} d (i (\xi) \omega). d \xi \right\} d t.
\]

We now deduce from formulas (15) and (16) the equalities

\[
\begin{array}{r l} & {d s (\omega) + s (d \omega) = \int_ {0} ^ {1} \left\{\int_ {\mathfrak {g}} \tau_ {\exp t \xi} ^ {*} (d (i (\xi) \omega) + i (\xi) d \omega). d \xi \right\} d t} \\ & {\qquad = \int_ {0} ^ {1} \left\{\int_ {\mathfrak {g}} \tau_ {\exp t \xi} ^ {*} (\theta (\xi) \omega). d \xi \right\} d t} \\ & {\qquad = \int_ {\mathfrak {g}} \left\{\int_ {0} ^ {1} \frac {d}{d t} (\tau_ {\exp t \xi} ^ {*} \omega) d t \right\} d \xi} \\ & {\qquad = \int_ {\mathfrak {g}} (\tau_ {\exp \xi} ^ {*} \omega - \omega) d \xi} \\ & {\qquad = \omega^ {\sharp} - \omega ,} \end{array}
\]

hence Th. 2.

We apply the theorem in the case X = G, for the action of G by left translations. Recall (Chap. III, §3, no. 13, Prop. 50) that associating to a differential form on G its value at the identity element gives an isomorphism from \(\Omega(\mathrm{G})^{\mathrm{G}}\) to the graded algebra \(\operatorname{Alt}(\mathfrak{g})\) of alternating multilinear forms on g. Identify \(\Omega(\mathrm{G})^{\mathrm{G}}\) with \(\operatorname{Alt}(\mathfrak{g})\) by means of this isomorphism. The operator d is then given by the formula (Chap. III, §3, no. 14, Prop. 51)

\[
\begin{array}{l} d \omega (a _ {1}, \ldots , a _ {p + 1}) \\ = \sum_ {i <   j} (- 1) ^ {i + j} \omega ([ a _ {i}, a _ {j} ], a _ {1}, \ldots , a _ {i - 1}, a _ {i + 1}, \ldots , a _ {j - 1}, a _ {j + 1}, \ldots , a _ {p + 1}) \end{array}
\]

for \(\omega\) in \(\mathrm{Alt}^p (\mathfrak{g})\) and \(a_1,\ldots ,a_{p + 1}\) in \(\mathfrak{g}\) .

For \(\xi \in \mathfrak{g}\) , let \(\mathrm{L}_{\xi}\) be the corresponding left-invariant vector field (defined by means of the action of G on itself by right translations, cf.~Chap. III, §3, no. 6). The operators \(\theta (\mathrm{L}_{\xi}), i(\mathrm{L}_{\xi})\) commute with the action of G on \(\Omega (\mathrm{G})\) defined by left translation, and hence induce operators \(\theta (\xi), i(\xi)\) on \(\Omega (\mathrm{G})^{\mathrm{G}}\) ; with the preceding identifications, these are expressed by the formulas (Differentiable and Analytic Manifolds, Results, 8.3.2 and 8.4.2)

\[
(\theta (\xi) \omega) (a _ {1}, \dots , a _ {p}) = - \sum_ {i} \omega (a _ {1}, \dots , a _ {i - 1}, [ \xi , a _ {i} ], a _ {i + 1}, \dots , a _ {p})
\]

\[
(i (\xi) \omega) (a _ {1}, \ldots , a _ {p - 1}) = \omega (\xi , a _ {1}, \ldots , a _ {p - 1})
\]

for \(\omega\) in \(\operatorname{Alt}^p (\mathfrak{g})\) and \(a_1,\ldots ,a_p\) in \(\mathfrak{g}\) .

The subcomplex \({}^{\mathrm{G}}\varOmega(\mathrm{G})^{\mathrm{G}}\) of biinvariant forms (Chap. III, §3, no. 13) can be identified with the subcomplex \(\operatorname{Alt}(\mathfrak{g})^{\mathrm{G}}\) of alternating multilinear forms on g invariant under the adjoint representation (that is, such that \(\theta(\xi)\omega=0\) for all \(\xi\in\mathfrak{g}\) ). Thus, we have a commutative diagram of complexes

\[
\begin{array}{c c c c c} ^ {\mathrm{G}} \varOmega (\mathrm{G}) ^ {\mathrm{G}} & \longrightarrow & \varOmega (\mathrm{G}) ^ {\mathrm{G}} & \longrightarrow & \varOmega (\mathrm{G}) \\ \Big \downarrow & & \Big \downarrow & & \\ \operatorname{Alt} (\mathfrak {g}) ^ {\mathrm{G}} & \longrightarrow & \operatorname{Alt} (\mathfrak {g}) & & \end{array}\tag{18}
\]

where the horizontal arrows are the canonical injections, and the vertical arrows are the isomorphisms induced by the map \(\omega\mapsto\omega(e)\) .

COROLLARY 1. a) In the diagram (18), all the morphisms are homotopisms. b) Let \(\omega \in \operatorname{Alt}(\mathfrak{g})\) . Then \(\omega\) belongs to \(\operatorname{Alt}(\mathfrak{g})^{\mathrm{G}}\) if and only if \(d\omega = 0\) and \(d(i(\xi)\omega) = 0\) for all \(\xi \in g\) . The differential of the complex \(\operatorname{Alt}(\mathfrak{g})^{\mathrm{G}}\) is zero. c) The graded vector space \(\mathrm{H}(\Omega(\mathrm{G}))\) is isomorphic to \(\operatorname{Alt}(\mathfrak{g})^{\mathrm{G}}\) .

The theorem, applied to the action of G on G by left translations (resp. to the action \(((g,h);x)\mapsto gxh^{-1}\) of G × G on G) implies that the canonical injection \(\Omega(\mathrm{G})^{\mathrm{G}}\to\Omega(\mathrm{G})\) (resp. \({}^{G}\Omega(\mathrm{G})^{\mathrm{G}}\to\Omega(\mathrm{G})\) ) is a homotopism; in view of Algèbre, Chap. X, p.~34, Cor., assertion a) follows.

We prove \(b)\) . By Prop. 51 of Chap. III, §3, no. 14, we have \(d\alpha = -d\alpha\) , that is \(d\alpha = 0\) , for every differential form \(\alpha\) on \(G\) that is left and right invariant. Thus, if \(\omega \in \mathrm{Alt}(\mathfrak{g})^G\) , then \(d\omega = 0\) , and consequently \(d(i(\xi)\omega) = \theta(\xi)\omega - i(\xi)d\omega = 0\) . Conversely, if \(d\omega = 0\) and \(d(i(\xi)d\omega) = 0\) , then \(\theta(\xi)\omega = 0\) . Assertion \(c)\) follows from \(a)\) and \(b)\) .

Remark. Consider the subcomplexes \(\mathrm{Z}(\varOmega(\mathrm{G}))\) and \(\mathrm{B}(\varOmega(\mathrm{G}))\) of \(\varOmega(\mathrm{G})\) (Algèbre, Chap. X, p.~25). It follows from the formula giving the differential of the product of two forms (Differentiable and Analytic Manifolds, Results, 8.3.5) that \(\mathrm{Z}(\varOmega(\mathrm{G}))\) is a subalgebra of \(\varOmega(\mathrm{G})\) and that \(\mathrm{B}(\varOmega(\mathrm{G}))\) is an ideal of \(\mathrm{Z}(\varOmega(\mathrm{G}))\) ; consequently, the exterior product induces a graded algebra structure on \(\mathrm{H}(\varOmega(\mathrm{G}))\) . The preceding now gives an isomorphism of graded algebras \(\mathrm{H}(\varOmega(\mathrm{G})) \to \mathrm{Alt}(\mathfrak{g})^{\mathrm{G}}\) .

Let H be a closed subgroup of G; we apply Th. 2 to X = G/H. By Chap. III, §1, no. 8, Cor. 1 of Prop. 17, the G-invariant differential forms on G/H can be identified with the H-invariant elements of \(\operatorname{Alt}(\mathrm{T}_{e}(\mathrm{G}/\mathrm{H}))\) , that is with the elements of \(\operatorname{Alt}(\mathfrak{g})\) that are H-invariant and annihilated by the operators \(i(\xi)\) for all \(\xi \in \mathrm{L}(\mathrm{H})\) . Consequently:

COROLLARY 2. Let \(\mathrm{H}\) be a closed subgroup of \(\mathrm{G}\) .

\begin{enumerate}
\def\labelenumi{\alph{enumi})}
\item
  The canonical injection \(\Omega (\mathrm{G / H})^{\mathrm{G}}\to \Omega (\mathrm{G / H})\) is a homotopism.
\item
  The complex \(\Omega(\mathrm{G}/\mathrm{H})^{\mathrm{G}}\) can be identified with the subcomplex of \(\operatorname{Alt}(\mathfrak{g})\) consisting of the elements \(\omega\) of \(\operatorname{Alt}(\mathfrak{g})\) that are invariant under the adjoint representation of H and are such that \(i(\xi)\omega = 0\) for all \(\xi \in \operatorname{L}(\operatorname{H})\) . If, in addition, H is connected, this subcomplex consists of the \(\omega \in \operatorname{Alt}(\mathfrak{g})\) such that \(\theta(\xi)\omega = 0\) and \(i(\xi)\omega = 0\) for all \(\xi \in \operatorname{L}(\operatorname{H})\) .
\end{enumerate}

